\documentclass[a4paper,11pt]{article}

\usepackage[english]{babel} 
\usepackage[utf8]{inputenc}
\usepackage[cyr]{aeguill}
\usepackage{stmaryrd}

\usepackage{lmodern} %Type1-font for non-english texts and characters
\usepackage{caption}
\usepackage{subcaption}

\usepackage{graphicx}
\usepackage{hyperref}

\usepackage{epstopdf}


%\hypersetup{	
%colorlinks=true,   %colorise les liens 
%breaklinks=true,  %permet le retour à la ligne dans les liens trop longs 
%urlcolor= blue,    %couleur des hyperliens 
%linkcolor= black, %couleur des liens internes 
%citecolor=black,	 %couleur des références 
%pdftitle={Compte rendu \emph{Traitement numérique du signal}, %informations apparaissant dans 
%pdfauthor={Mélisande Zonta}, %les informations du document 
%pdfsubject={Projet TNS2}	%sous Acrobat. 
%} 

%% Math Packages 
\usepackage{amsmath}
\usepackage{amsthm}
\usepackage{amsfonts}
\usepackage{amssymb}
\usepackage{mathrsfs}
\usepackage{pst-all}
\usepackage{lscape}
\usepackage{pdfpages}
\usepackage{mathabx}


\usepackage{color, colortbl}
\definecolor{lightgray}{gray}{0.85}
\usepackage{multirow}
\usepackage[Algorithme]{algorithm}
\usepackage[noend]{algpseudocode}
\usepackage{tikz}

%\renewcommand{\algorithmicdo}{\textbf{faire}}
%\renewcommand{\algorithmicwhile}{\textbf{tant que}}


\usepackage{a4wide} %%Smaller margins = more text per page.
\usepackage{fancyhdr} %%Fancy headings

\setcounter{secnumdepth}{5}
\setcounter{tocdepth}{5}


\DeclareMathOperator*{\argmax}{arg\,max}
\DeclareMathOperator*{\argmin}{arg\,min}
\graphicspath{{/Users/melisandezonta/Documents/Documents/GTL_courses_second_semester/Computer-Vision/PS2-all/PS2-images/}}

\begin{document}

%\pagestyle{fancy}

\begin{titlepage}
\vspace*{\stretch{1}}

\begin{center}
\includegraphics[scale=0.4]{GT_logo.jpeg}
\end{center}
\vspace*{\stretch{1}}
\hrulefill
\begin{center}\bfseries\huge
   Computer Vision \\
   CS 6476 , Spring 2018\\
   \end{center}
  \begin{center}\bfseries\large
     PS 2\\
    \hrulefill
\end{center}
%\hfill
\vspace*{1cm}
\begin{minipage}[t]{0.6\textwidth}
  \begin{flushleft} \large
    \emph{Professor : }\\
    Cedric Pradalier \\
  \end{flushleft}
\end{minipage}
\begin{minipage}[t]{0.3\textwidth}
  \begin{flushright} \large
    \emph{Author :} \\
    Melisande Zonta \\
  \end{flushright}
\end{minipage}
\vspace*{\stretch{2}}
\begin{flushright}
       \today 
\end{flushright} 
\end{titlepage}

\tableofcontents
\clearpage



\section{Implementation of the SSD match algorithm}

The Figure 1 a. and b. show respectively the left to right and right to left disparity results using the SSD scoring method.

 \begin{figure}[H]
\begin{center}
\begin{tabular}{cc}
	\includegraphics[width=.5\textwidth]{ps2-1-a-lefttoright.png}&
	\includegraphics[width=.5\textwidth]{ps2-1-a-righttoleft.png}\\
	a&b
\end{tabular}
\end{center}
\caption{ 
\textit{a}. ps2-1-a : Left to right Image.  \textit{b}. ps2-1-a : Right to left Image. }
\label{ps2-1}
\end{figure}

\section{Test on original images}

\subsection{Application of the SSD algorithm}



 \begin{figure}[H]
\begin{center}
\begin{tabular}{cc}
	\includegraphics[width=.5\textwidth]{ps2-2-a-lefttoright.png}&
	\includegraphics[width=.5\textwidth]{ps2-2-a-righttoleft.png}\\
	a&b
\end{tabular}
\end{center}
\caption{ 
\textit{a}. ps2-2-a : Left to right Image.  \textit{b}. ps2-2-a : Right to left Image. }
\label{ps2-1}
\end{figure}

The Figure 2 a. and b.show the result of the same algorithm, this time applied to the real picture of the artist?s workshop.

The images were obtained for a window size of 9 pixels and a maximum measured disparity of 100 pixels.

First we notice that our results for the left to right and right to left disparity images are consistent. We more or less get the same output, with a darker region on one of the side, which is due to the fact that we cannot compute the disparity on one edge, because the corresponding pixels actually lie outside our target image.

What?s more, we can say that most portions of our disparity maps match the ground truth, especially for large objects. Thinner objects such as the pens and paintbrushes? handles do not get matched. The main reason is that the window in this case covers areas located at varying depth, which makes it difficult to find a match. The larger brush, on the left, is better covered.

The same issue arises regarding the vertical bars supporting the red rings in the back, but with an additional issue due to the repeating background which makes it difficult to distinguish the individual bars.

The rest of the objects, including the ring, seem to get a reasonable disparity value, which is why we decided to keep this window size that gives pretty good results.

\subsection{Comparison of the results with the ground truth}

 \begin{figure}[H]
\begin{center}
\begin{tabular}{cc}
	\includegraphics[width=.5\textwidth]{ps2-2-b-ground_truth_left.png}&
	\includegraphics[width=.5\textwidth]{ps2-2-b-ground_truth_right.png}\\
	a&b
\end{tabular}
\end{center}
\caption{ 
\textit{a}. ps2-2-b : Ground Truth Left Image.  \textit{b}.  ps2-2-b : Ground Truth Right Image. }
\label{ps2-1}
\end{figure}

\section{Robustness of the SSD algorithm to certain perturbations}

\subsection{Add of Gaussian noise of noise}


 \begin{figure}[H]
\begin{center}
\begin{tabular}{cc}
	\includegraphics[width=.5\textwidth]{ps2-3-a-lefttoright.png}&
	\includegraphics[width=.5\textwidth]{ps2-3-a-righttoleft.png}\\
	a&b
\end{tabular}
\end{center}
\caption{ 
\textit{a}. ps2-3-a : Left to right Image.  \textit{b}. ps2-3-a : Right to left Image. }
\label{ps2-1}
\end{figure}

Those results (Figure 3 a. and b.) were obtained for a window size of 9 pixels and a noise sigma of 25.

We use a quite high sigma value in order to get a very clear idea of the effect of a Gaussian noise on the disparity computation.
Using the same parameters as in the previous question, we notice that we can still clearly see the similarities between the two disparities, but the erroneous pixels created a lot of disparity noise.

Regions in large objects that appeared correctly with a rather smooth disparity now have a more granular aspect. The erroneous patches already visible in the previous test are still visible, with additional smaller errors spread across the disparity map.

However, we notice that although this as a big impact in regions with smooth disparity (based on the ground truth), the other regions of the image, which were not properly matched in the previous test, are not as much affected by the noise (example of the pens and vertical bars).


\subsection{Increasing of the contrast}



 \begin{figure}[H]
\begin{center}
\begin{tabular}{cc}
	\includegraphics[width=.5\textwidth]{ps2-3-b-lefttoright.png}&
	\includegraphics[width=.5\textwidth]{ps2-3-b-righttoleft.png}\\
	a&b
\end{tabular}
\end{center}
\caption{ 
\textit{a}. ps2-3-b : Left to right Image.  \textit{b}. ps2-3-b : Right to left Image. }
\label{ps2-1}
\end{figure}


In comparison, increasing the contrast simply by multiplying one of the image by a scaling factor has a very different impact on the output of the disparity computation. The result is visible in Figure 4 a. and b. 

Most regions of smooth disparity that we had in the first test now get ?patchier?, they get broken up in smaller smoothed patches of disparity.

Although the "edges" of the objects seem to be conserved, this looks like a worse result than in the Gaussian noise test.
Indeed, whereas some smoothing on the disparity map could probably solve some problems caused by the noise, this is trickier with the contrast effect, as the mean values of larger regions has been lost.

\section{Implementation of the normalized correlation}

\subsection{Application of the algorithm}

The tests below were done using a disparity computation based on the normalized cross-correlation available in through the cv2.matchTemplate function.

 \begin{figure}[H]
\begin{center}
\begin{tabular}{cc}
	\includegraphics[width=.5\textwidth]{ps2-4-a-lefttoright.png}&
	\includegraphics[width=.5\textwidth]{ps2-4-a-righttoleft.png}\\
	a&b
\end{tabular}
\end{center}
\caption{ 
\textit{a}. ps2-4-a : Left to right Image.  \textit{b}. ps2-4-a : Right to left Image. }
\label{ps2-1}
\end{figure}


Without noise (Figure 5 a. and b.), it seems that large areas of similar disparity get correctly matched but we have weird effects appearing as ?shadows? casted in the opposite direction of the matching.

\subsection{Robustness Tests}

 \begin{figure}[H]
\begin{center}
\begin{tabular}{cc}
	\includegraphics[width=.5\textwidth]{ps2-4-b-noise-lefttoright.png}&
	\includegraphics[width=.5\textwidth]{ps2-4-b-noise-righttoleft.png}\\
	a&b
\end{tabular}
\end{center}
\caption{ 
\textit{a}. ps2-4-b : Left to right Image.  \textit{b}. ps2-4-b : Right to left Image. }
\label{ps2-1}
\end{figure}



 \begin{figure}[H]
\begin{center}
\begin{tabular}{cc}
	\includegraphics[width=.5\textwidth]{ps2-4-b-noise-lefttoright.png}&
	\includegraphics[width=.5\textwidth]{ps2-4-b-ground_truth_right.png}\\
	a&b
\end{tabular}
\end{center}
\caption{ 
\textit{a}. ps2-4-b : Left to right Image.  \textit{b}. ps2-4-b : Right to left Image. }
\label{ps2-1}
\end{figure}


 \begin{figure}[H]
\begin{center}
\begin{tabular}{cc}
	\includegraphics[width=.5\textwidth]{ps2-4-b-ground_truth_left.png}&
	\includegraphics[width=.5\textwidth]{ps2-4-b-ground_truth_right.png}\\
	a&b
\end{tabular}
\end{center}
\caption{ 
\textit{a}. ps2-4-b : Ground Truth Left Image.  \textit{b}. ps2-4-b : Ground Truth Right Image. }
\label{ps2-1}
\end{figure}


The results from images subjected to Gaussian noise (Figure 6 a. and b.) are suspicious.

The matching does not really work, and an ?edge effect? is visible on the disparity map, with only the boundaries of large objects distinguishable.

Images resulting from increased contrast versions of the original one (Figure 7 a and b.) look very similar to the one obtained from the original images.

It seems that  the cross-correlation method is not affected by contrast changes as the SSD one.

As a conclusion for this part, we could say that the cross-correlation method is more efficient than the SSD method.

Indeed, taking a proper implementation of this method, the fact that it does not suffer from contrast changes seems quite positive for real case stereo matching scenarios, where left and right images could slightly differ in contrast.

Furthermore, the apparent sensitivity of this method to Gaussian noise is not a big issue in modern applications, where this type of noise appears mainly in situations where the lights are low, typically at night.

If the setup is more favorable, this should not be a problem, especially considering that our tests were done with rather high level of noise, which is quite improbable in real case scenarios.

\section{Final Test on other images}


 \begin{figure}[H]
\begin{center}
\begin{tabular}{cc}
	\includegraphics[width=.5\textwidth]{ps2-5-original-ssd-lefttoright.png}&
	\includegraphics[width=.5\textwidth]{ps2-5-original-ssd-righttoleft.png}\\
	a&b
\end{tabular}
\end{center}
\caption{ 
\textit{a}. ps2-5-ssd : Original left to right image.  \textit{b}. ps2-5-ssd : Original right to left image. }
\label{ps2-1}
\end{figure}


 \begin{figure}[H]
\begin{center}
\begin{tabular}{cc}
	\includegraphics[width=.5\textwidth]{ps2-5-original-corr-lefttoright.png}&
	\includegraphics[width=.5\textwidth]{ps2-5-original-corr-lefttoright.png}\\
	a&b
\end{tabular}
\end{center}
\caption{ 
\textit{a}. ps2-5-corr : Original left to right image.  \textit{b}. ps2-5-corr : Original right to left image. }
\label{ps2-1}
\end{figure}

% Ground truth

 \begin{figure}[H]
\begin{center}
\begin{tabular}{cc}
	\includegraphics[width=.5\textwidth]{ps2-5-ground_truth_left.png}&
	\includegraphics[width=.5\textwidth]{ps2-5-ground_truth_right.png}\\
	a&b
\end{tabular}
\end{center}
\caption{ 
\textit{a}. ps2-5-ground-truth : Left Image.  \textit{b}. ps2-5-ground-truth : Right Image. }
\label{ps2-1}
\end{figure}


% Blur or noise


\begin{figure}[H]
\begin{center}
\begin{tabular}{cc}
	\includegraphics[width=.5\textwidth]{ps2-5-noise-ssd-lefttoright.png}&
	\includegraphics[width=.5\textwidth]{ps2-5-noise-ssd-righttoleft.png}\\
	a&b
\end{tabular}
\end{center}
\caption{ 
\textit{a}. ps2-5-ssd-blur : Left to right image after blurring.  \textit{b}. ps2-5-ssd-blur: Right to left image after blurring. }
\label{ps2-1}
\end{figure}


 \begin{figure}[H]
\begin{center}
\begin{tabular}{cc}
	\includegraphics[width=.5\textwidth]{ps2-5-noise-corr-lefttoright.png}&
	\includegraphics[width=.5\textwidth]{ps2-5-noise-corr-righttoleft.png}\\
	a&b
\end{tabular}
\end{center}
\caption{ 
\textit{a}. ps2-5-corr-blur : Left to right image after blurring.  \textit{b}. ps2-5-corr-blur: Right to left image after blurring. }
\label{ps2-1}
\end{figure}

% Contrast on both images

\begin{figure}[H]
\begin{center}
\begin{tabular}{cc}
	\includegraphics[width=.5\textwidth]{ps2-5-contrast-ssd-lefttoright.png}&
	\includegraphics[width=.5\textwidth]{ps2-5-contrast-ssd-righttoleft.png}\\
	a&b
\end{tabular}
\end{center}
\caption{ 
\textit{a}. ps2-5-ssd-contrast : Left to right image after blurring.  \textit{b}. ps2-5-ssd-contrast: Right to left image after blurring. }
\label{ps2-1}
\end{figure}


 \begin{figure}[H]
\begin{center}
\begin{tabular}{cc}
	\includegraphics[width=.5\textwidth]{ps2-5-contrast-corr-lefttoright.png}&
	\includegraphics[width=.5\textwidth]{ps2-5-contrast-corr-lefttoright.png}\\
	a&b
\end{tabular}
\end{center}
\caption{ 
\textit{a}. ps2-5-corr-contrast : Left to right image after blurring.  \textit{b}. ps2-5-corr-contrast: Right to left image after blurring. }
\label{ps2-1}
\end{figure}


\end{document}